\section{AMI Labs: emprender junto con LeCun}

Xie Saining describe el paso de la investigación al emprendimiento como una curva de extensión natural, no como un giro repentino. Para él, world model no es una palabra que persigue una tendencia, sino el resultado de que representation learning siga avanzando. Lo que en verdad cambió el ritmo fue esa conversación con Yann LeCun en 2025. En un principio fue su mentor quien le sugirió que bien podía preguntarle a LeCun si estaría dispuesto a hacer algo juntos, pero en la reunión uno a uno de la segunda semana fue el propio LeCun quien propuso emprender. Las visiones de ambos coincidían en gran medida: los dos buscaban building the predictive brain. Xie Saining describe esa sensación así: \textit{«hablar con Yann es un poco como casting spells... dice algo y ya no puedes pensar en nada más»}, y solo tardó una semana en decidir sumarse.

\begin{importantbox}{El punto de partida de AMI Labs: la visión antes que la organización}
Lo más interesante de AMI Labs es que no partió de un caparazón corporativo al que después se le buscara una historia capaz de atraer financiamiento, sino de un alto consenso previo sobre la dirección de investigación «predictive brain», a partir del cual se organizaron el equipo, el capital y la estructura de ejecución. La razón por la que Xie Saining se sumó tampoco fue que «emprender resulte más rentable», sino haber visto una ventana poco común: la posibilidad de llevar el modelo del mundo de una tesis académica a una ingeniería organizacional orientada a un producto real y a una pila tecnológica de largo plazo.
\end{importantbox}

Este camino resulta válido también porque Xie Saining ya había rechazado en dos ocasiones anteriores oportunidades más prestigiosas, y que el exterior consideraría con mayor facilidad como la opción «correcta». La primera vez fue al graduarse del doctorado, cuando rechazó a OpenAI, lo que incluso llevó a Ilya Sutskever a llamarlo por teléfono para expresarle su fuerte desacuerdo; la segunda fue en 2024, cuando volvió a rechazar a Ilya en la fundación de SSI. Si se observan juntos estos dos rechazos y la elección de AMI, se advierte que Xie Saining no se limitaba a «alejarse de la industria», sino que buscaba de manera constante un espacio más acorde con su propia manera de plantear los problemas. Si una organización solo trata el modelo del mundo como un empaque de PR, o solo permite optimizar en torno a métricas de producto a corto plazo, entonces, por más recursos que tenga, no necesariamente es adecuada para él.

Esto explica por qué afirma con claridad que no quiere sumarse a las grandes empresas actuales de Silicon Valley ni a un lab cerrado. Según su observación, la tónica dominante en Silicon Valley en los últimos años ha sido el cierre, el secreto, no publicar artículos y no abrir el código fuente, lo cual genera, en cambio, un rechazo natural hacia la investigación exploratoria. La carrera armamentista de las grandes empresas ha exprimido el oxígeno de research: muchos equipos, aunque en teoría conservan una función de investigación, en la práctica quedan por completo arrastrados por los cronogramas de producto y la disputa por recursos de cómputo. Lo que Xie Saining busca no es un research lab puro ni tampoco una empresa de grandes modelos completamente cerrada, sino un punto de equilibrio muy poco común entre ambos extremos: lo bastante serio para funcionar como una empresa emergente, capaz de sostener organizacionalmente una inversión de largo plazo, y a la vez con suficiente research freedom para no comprimir todo el trabajo en KPI trimestrales.

\begin{warningbox}{«Ir al lugar con más dinero para hacer la investigación más vanguardista» ya no se cumple de manera automática}
La crítica de Xie Saining hacia Silicon Valley no es un simple sentimiento territorial, sino un juicio sobre el cambio en el entorno de investigación. Cuando el secreto, la carrera armamentista y las expectativas de capital se imponen sobre la discusión abierta, el lugar con más recursos no necesariamente es el más adecuado para el trabajo más exploratorio. Confundir «más dinero» con «mejor entorno de investigación» es un error cognitivo en el que hoy caen con facilidad no pocos investigadores jóvenes.
\end{warningbox}

Por eso el diseño organizacional de AMI Labs también resulta bastante simbólico. El equipo inicial cuenta con alrededor de 25 personas, una valuación de aproximadamente 3,000 millones de dólares, e incluye a 6 co-founder, un CEO y al directivo principal a cargo de world model. Más impresionante aún es que varios cofundadores renunciaron a decenas de millones de dólares en stock ya vested de OpenAI o Meta para volver a sentarse a la mesa. Este detalle muestra que los miembros del equipo no perseguían el simple juego financiero de «una valuación más alta», sino que apostaban por una dirección tecnológica que, a su juicio, se acerca más a la línea principal del futuro que las plataformas ya existentes. Para Xie Saining, esto también coincide con su criterio de juicio habitual: lo importante no es qué compañía parece hoy más cercana al desenlace final, sino con quién y en qué estás trabajando en este momento.

\begin{knowledgebox}{Por qué Nueva York, y no Silicon Valley}
AMI Labs eligió Nueva York tanto por la larga trayectoria académica de Yann LeCun en NYU como por un juicio deliberado sobre el carácter de la ciudad. La frase de Xie Saining, \textit{``Silicon Valley is very unpeeled''}, tiene un matiz muy personal y viene a decir que ese lugar se deja hipnotizar con demasiada facilidad por el lujo, las valuaciones y la ilusión de la cámara de eco. Nueva York, en cambio, se parece más a un corte transversal del mundo real: las finanzas, el arte, la migración, la calle, el cine y la academia coexisten ahí, de modo que uno no llega a confundir la IA con el mundo mismo.
\end{knowledgebox}

En la parte sobre la experiencia de emprender hay también una metáfora muy cargada de sentido de acción: esquiar. Xie Saining dice que, al esquiar, hay que orientar los hombros pendiente abajo sin ningún temor, algo muy contrario al instinto, porque cuando una persona tiene miedo tiende de forma instintiva a echarse hacia atrás, frenar y tratar de despegarse de la pendiente; pero es precisamente ese movimiento defensivo el que con mayor facilidad provoca una pérdida de control real. Con emprender ocurre lo mismo. Cuando los problemas se vuelven cada vez más grandes y las responsabilidades cada vez más concretas, no basta con pensar en conservar la posición cómoda que se tenía dentro del sistema anterior; hay que girar el cuerpo de manera activa hacia la dirección más empinada, que es, a la vez, la única que en verdad lleva hacia adelante.

\begin{importantbox}{Emprender se parece a esquiar, y también a investigar}
Esquiar, investigar y emprender comparten una estructura muy profunda: en un entorno de alta incertidumbre, la reacción más natural de una persona suele ser encogerse, pero la acción realmente eficaz suele consistir en salir al encuentro del problema de manera más activa. Con esta metáfora, Xie Saining explica que lo que valora en un miembro del equipo no es solo un currículum vistoso, sino si de verdad tiene el valor de orientar los hombros pendiente abajo, y si está dispuesto a enfrentar durante mucho tiempo un problema difícil, en lugar de parecer comprometido únicamente cuando las cosas van a favor.
\end{importantbox}

En la entrevista hay también un punto muy importante y que se pasa por alto con facilidad: AMI Labs no se presenta a sí misma como una empresa emergente «antiacadémica». Al contrario, lo que intenta heredar es esa integridad de científico propia de LeCun, y al mismo tiempo reconoce que, si el modelo del mundo de verdad importa, no puede quedarse para siempre en los paper y los talk. Investigación y emprendimiento no son aquí dos bandos opuestos, sino dos maneras de hacer avanzar el mismo problema en escalas de tiempo distintas: la primera se encarga de hacer concreto el concepto, y la segunda, de situar a la organización, el capital y el talento en una estructura capaz de sostener los objetivos de largo plazo.

En este sentido, la decisión de Xie Saining de emprender resulta, en realidad, bastante coherente. No pasó de ser un «idealista de la investigación» a un «realista de los negocios», sino que colocó la línea de visión y modelo del mundo en la que había creído durante mucho tiempo dentro de un nuevo terreno que exige asumir consecuencias mayores. Aquí, desde luego, el riesgo es más alto y el ritmo más rápido, pero, tal como muestra una y otra vez en la entrevista, no le teme a colocarse dentro de un sistema de alta incertidumbre, siempre que la manera de plantear los problemas de ese sistema sea suficientemente real.

\subsection{Resumen de la sección}

La aparición de AMI Labs es a la vez una extensión natural de la investigación de Xie Saining sobre el modelo del mundo y una práctica más de su rechazo constante a tomar el brillo de una plataforma a corto plazo como punto final. Haber rechazado a Ilya en dos ocasiones, rechazar la lógica cerrada de Silicon Valley y elegir Nueva York para emprender junto con LeCun muestra que lo que en verdad persigue es un entorno organizacional capaz de construirse tomando predictive brain como línea principal.
